% Lliçó: límit d'una funció en un punt (valencià).
%
% El mateix contingut que es.tex, frase a frase i fórmula a fórmula. Els noms
% dels entorns (Definició, Nota, Demostració...) no s'escriuen: els posa
% l'idioma en què es compila.

\begin{frame}
\didactatitle{Límit d'una funció en un punt}

\bymedium{%
  $f(x)$ s'acosta a $L$ quan $x$ s'acosta a $a$. La definició diu què vol
  dir «s'acosta».
}{%
  La idea intuïtiva és que $f(x)$ s'acosta tant com es vulga a $L$ quan $x$
  s'acosta prou a $a$. La definició converteix aquesta frase en una
  desigualtat que es pot comprovar.
}

\begin{definition}[Límit]
Siga $f$ una funció definida en un interval obert que conté $a$, excepte
potser en el mateix $a$. Diem que $f$ té \keyterm{límit} $L$ en $a$, i
escrivim $\lim_{x \to a} f(x) = L$, si per a cada $\varepsilon > 0$ existeix
un $\delta > 0$ tal que
\begin{keyformula}
  0 < |x - a| < \delta \implies |f(x) - L| < \varepsilon .
\end{keyformula}
\end{definition}

\begin{notesonly}
La condició $0 < |x - a|$ deixa fora el mateix punt $a$: el límit no depén
del valor de $f$ en $a$, ni tan sols de si hi està definida. Només mira què
passa \emph{a prop} de $a$. Per això té sentit preguntar-se pel límit de
$(x^2 - 1)/(x - 1)$ en $x = 1$, on el quocient no està definit.

A més, el límit, si existeix, és únic: dos valors diferents $L \neq L'$ no
poden estar alhora a menys de $\varepsilon = |L - L'|/2$ de $f(x)$.
\end{notesonly}

\begin{teaching}[Ritme]
Vint minuts. Dibuixa la banda $(L - \varepsilon, L + \varepsilon)$ abans
d'escriure la fórmula: l'ordre dels quantificadors (primer et donen
$\varepsilon$, després tries $\delta$) s'entén millor a la pissarra.
\end{teaching}
\end{frame}
